\chapter{Составные системы, запутанность и синглет}

\lectureoverview{Вопрос об испускании фотона выявляет различие между квантовым событием и его средним по ансамблю. Это различие естественно приводит к измерению, коллапсу и необходимости рассматривать систему вместе с её аппаратом. Поэтому мы строим пространство состояний тензорного произведения двух систем, затем специализируемся на двух спинах, отличаем продуктные состояния от запутанных и встраиваем каждую односпиновую наблюдаемую в составное пространство. Затем синглетное состояние даёт решающий пример: каждый спин локально неполяризован, однако одинаковые компоненты идеально антикоррелированы. Его базис-независимая форма подготавливает почву для теоремы Белла и для различения неклассической корреляции и передачи сигналов быстрее света.}

\section{Дискретные переходы для спина в суперпозиции}

Начнём с одного спина в магнитном поле вдоль оси \(z\). В единицах, где \(\hbar=1\),
\begin{equation}
H=\frac{\omega}{2}\sigma_z,
\qquad
E_+=+\frac{\omega}{2},
\qquad
E_-=-\frac{\omega}{2}.
\label{eq:l6-spin-hamiltonian}
\end{equation}
Два уровня разделены на
\begin{equation}
\Delta E=E_+-E_-=\omega.
\label{eq:l6-energy-gap}
\end{equation}
При учёте поля излучения спин на верхнем уровне может перейти на нижний уровень и испустить фотон энергии \(\omega\). У спина, уже находящегося на нижнем уровне, нет более низкого спинового состояния, в которое он мог бы излучательно перейти.

\begin{figure}[htbp]
\centering
\includegraphics[width=0.86\textwidth]{lecture_06_figure_01.png}
\caption{Два спиновых энергетических уровня и их энергетический зазор \(\omega\) при \(\hbar=1\).}
\lectureframe{00:02:35}
\end{figure}

Тонкий случай — это спин, подготовленный поперёк поля. У классического магнитного ротора энергия меняется непрерывно и пропорциональна косинусу угла между ним и полем. Он прецессирует с угловой частотой \(\omega\), излучает на этой же частоте и постепенно опускается к ориентации минимальной энергии. Полная излучённая энергия непрерывно зависит от начального угла: ротор, изначально горизонтальный, выделяет вдвое меньше энергии, чем ротор, который вначале находится наверху.

Квантовая механика сохраняет этот классический ответ только в среднем. Запишем
\begin{equation}
\ket{\psi}=\alpha\ket{u}+\beta\ket{d},
\qquad
|\alpha|^2+|\beta|^2=1,
\label{eq:l6-superposed-spin}
\end{equation}
где \(\ket{u}\) и \(\ket{d}\) — верхнее и нижнее собственные состояния энергии. Верхняя компонента может распасться; нижняя не может. Поэтому
\begin{equation}
P(\text{one photon})=|\alpha|^2,
\qquad
P(\text{no photon})=|\beta|^2.
\end{equation}
Если фотон испускается, его энергия равна полному межуровневому зазору \(\omega\), а не его доле. Тем не менее средняя испущенная энергия равна
\begin{equation}
\expect{E_{\rm emitted}}=|\alpha|^2\omega.
\end{equation}
Для спина вдоль \(x\) имеем \(|\alpha|^2=|\beta|^2=1/2\), так что среднее равно \(\omega/2\). Каждое отдельное событие по-прежнему даёт либо один фотон энергии \(\omega\), либо ни одного фотона. Непрерывный классический результат относится к среднему по ансамблю; дискретность относится к каждому событию.

\begin{classroomqa}
\audiencequestion{Насколько точно энергия поглощаемого фотона должна совпадать с межуровневым зазором?}
\lecturerresponse{Ширина резонанса конечна. Если время жизни возбуждённого состояния равно \(\tau\), его характерная энергетическая ширина имеет порядок \(\Delta E\sim\hbar/\tau\). Больший магнитный дипольный момент обычно сильнее связан с полем излучения, сокращает время жизни и уширяет линию. Это соотношение неопределённостей энергия–время в форме, существенной для неустойчивых уровней.}
\end{classroomqa}

\begin{classroomqa}
\audiencequestion{Может ли само испускание фотона служить аппаратом, измеряющим спин?}
\lecturerresponse{Оно может показать, произошёл ли распад, но это не та идеальная повторяемая модель измерения, которая здесь нужна. Измерение компоненты спина должно оставлять спин в собственном состоянии, связанном с зарегистрированным значением. Напротив, распад переводит каждую ветвь с испусканием на нижний уровень, так что он не просто регистрирует оба возможных входных значения, сохраняя соответствующие им собственные состояния.}
\end{classroomqa}

\section{Коллапс и необходимость составной системы}

Между вмешательствами изолированное состояние эволюционирует по уравнению Schr\"odinger. Измерение обычно формулируют как дополнительное правило. Пусть \(L\) имеет ортонормированные собственные векторы \(\ket{\lambda_i}\), и разложим
\begin{equation}
\ket{\psi}=\sum_i\alpha_i\ket{\lambda_i},
\qquad
L\ket{\lambda_i}=\lambda_i\ket{\lambda_i}.
\label{eq:l6-measurement-expansion}
\end{equation}
Измерение даёт одно значение \(\lambda_k\) с вероятностью \(|\alpha_k|^2\), после чего состояние имеет вид
\begin{equation}
\ket{\psi}\longrightarrow\ket{\lambda_k}.
\label{eq:l6-collapse}
\end{equation}
Исчезновение остальных компонент из описания после измерения называется коллапсом.

\begin{figure}[htbp]
\centering
\includegraphics[width=0.84\textwidth]{lecture_06_figure_02.png}
\caption{Состояние, разложенное по собственному базису наблюдаемой \(L\), до того как выбран один из исходов.}
\lectureframe{00:12:35}
\end{figure}

\begin{classroomqa}
\audiencequestion{Что не так с правилом коллапса?}
\lecturerresponse{С ним всё в порядке как с операционным правилом для предсказания результатов измерений. Концептуальная трудность состоит в том, что аппарат сам является физической системой. Если и объект, и аппарат подчиняются квантовой механике, почему их взаимодействие должно требовать второго динамического закона в дополнение к унитарной эволюции Schr\"odinger? Здесь мы пойдём по пути описания объекта и аппарата как одной большей квантовой системы.}
\end{classroomqa}

Жёсткая граница между системой и аппаратом была центральным пунктом исторических споров. Вместо того чтобы решать этот спор декларативно, мы сначала разберёмся, как объединяются квантовые системы. Как только пространство состояний объединённой системы построено, измерение можно изучать как взаимодействие, коррелирующее систему с аппаратом. Поэтому ближайший математический вопрос таков:
\begin{quote}
Если даны две квантовые системы, каково пространство состояний целого?
\end{quote}

\section{Пространство состояний тензорного произведения}

Обозначим системы как \(A\) и \(I\), с гильбертовыми пространствами \(\mathcal S_A\) и \(\mathcal S_I\). Их составное пространство равно
\begin{equation}
\mathcal S_{AI}=\mathcal S_A\otimes\mathcal S_I.
\label{eq:l6-tensor-space}
\end{equation}
Выберем ортонормированные базисы \(\{\ket{a}\}\) и \(\{\ket{i}\}\) для двух множителей. Сопоставим каждому базисному вектору первой системы каждый базисный вектор второй:
\begin{equation}
\ket{a}\otimes\ket{i}\equiv\ket{a\,i}.
\end{equation}
Если пространства-множители имеют размерности \(D_A\) и \(D_I\), то
\begin{equation}
\dim\mathcal S_{AI}=D_A D_I.
\end{equation}
Скалярное произведение факторизуется,
\begin{equation}
\braket{b\,j}{a\,i}
=\braket{b}{a}\braket{j}{i}
=\delta_{ba}\delta_{ji}.
\label{eq:l6-tensor-inner-product}
\end{equation}
Следовательно, два составных базисных состояния ортогональны всякий раз, когда различается хотя бы одна метка подсистемы.

\begin{figure}[htbp]
\centering
\includegraphics[width=0.86\textwidth]{lecture_06_figure_03.png}
\caption{Отдельные пространства состояний, их тензорное произведение и парный ортонормированный базис.}
\lectureframe{00:24:05}
\end{figure}

Парные базисные векторы не образуют всего пространства. Их полная линейная оболочка имеет вид:
\begin{equation}
\ket{\Psi}=\sum_{a,i}c_{ai}\ket{a\,i},
\qquad
\sum_{a,i}|c_{ai}|^2=1.
\label{eq:l6-general-tensor-state}
\end{equation}

\begin{classroomqa}
\audiencequestion{Являются ли \(\ket{a}\) и \(\ket{i}\) произвольными векторами?}
\lecturerresponse{Здесь это элементы выбранных ортонормированных базисов. В любом из множителей можно использовать любой ортонормированный базис. Попарное объединение этих базисных векторов строит ортонормированный базис составного пространства, после чего произвольные векторы получаются линейными комбинациями, как в уравнении~\eqref{eq:l6-general-tensor-state}.}
\end{classroomqa}

\begin{classroomqa}
\audiencequestion{Зависит ли тензорное произведение от выбранных базисов?}
\lecturerresponse{Его координатное описание зависит, но само пространство — нет. Замена базиса в любой подсистеме лишь меняет индуцированный базис, используемый для описания тех же самых векторов в \(\mathcal S_A\otimes\mathcal S_I\).}
\end{classroomqa}

\section{Два спина и продуктные состояния}

Теперь рассмотрим два спина. Компоненты первого обозначим через \(\sigma_x,\sigma_y,\sigma_z\), а компоненты второго — через \(\tau_x,\tau_y,\tau_z\). В \(z\)-базисе составной базис имеет вид
\begin{equation}
\ket{uu},\qquad \ket{ud},\qquad \ket{du},\qquad \ket{dd}.
\label{eq:l6-two-spin-basis}
\end{equation}
Первая позиция всегда относится к \(\sigma\), вторая — к \(\tau\). Общее состояние имеет вид
\begin{equation}
\ket{\Psi}
=c_{uu}\ket{uu}+c_{ud}\ket{ud}
+c_{du}\ket{du}+c_{dd}\ket{dd}.
\label{eq:l6-general-two-spin}
\end{equation}

\begin{classroomqa}
\audiencequestion{Если два спина взаимодействуют, меняются ли сами базисные векторы?}
\lecturerresponse{Чтобы прослеживать динамику, не требуется никакого фиксированного базиса. В выбранном базисе взаимодействие изменяет коэффициенты \(c_{ab}(t)\), а не определение \(\ket{uu},\ket{ud},\ket{du},\ket{dd}\). Состояние, которое начинается как \(\ket{uu}\), может эволюционировать в суперпозицию или даже в \(\ket{dd}\), тогда как координатный базис остаётся тем же.}
\end{classroomqa}

Эта конструкция должна также пройти проверку на согласованность. Если спины далеко разнесены и рассматривается только \(\sigma\), предсказания должны согласовываться с односпиновой теорией, разработанной до введения \(\tau\). Тензорные произведения специально устроены так, чтобы сохранять эту локальную теорию.

Некоторые двухспиновые состояния описывают независимо подготовленные спины. Например, пусть
\begin{equation}
\ket{R}=\frac{\ket{u}+\ket{d}}{\sqrt2}
\end{equation}
будет собственным состоянием компоненты по \(x\) с собственным значением \(+1\). Тогда
\begin{equation}
\ket{u}_{\sigma}\otimes\ket{R}_{\tau}
=\frac{\ket{uu}+\ket{ud}}{\sqrt2}.
\label{eq:l6-up-right-product}
\end{equation}

\begin{classroomqa}
\audiencequestion{Можно ли подготовить два множителя вдоль разных осей?}
\lecturerresponse{Да. Каждый множитель может быть любым односпиновым состоянием. В уравнении~\eqref{eq:l6-up-right-product} первый спин определён вдоль \(z\), а второй — вдоль \(x\). Разложение второго множителя по \(z\)-базису лишь переписывает то же самое продуктное состояние в составном базисе.}
\end{classroomqa}

Наиболее общее продуктное состояние имеет вид
\begin{align}
\ket{\Psi_{\rm prod}}
&=(\alpha_u\ket{u}+\alpha_d\ket{d})
\otimes(\beta_u\ket{u}+\beta_d\ket{d})\\
&=\alpha_u\beta_u\ket{uu}
+\alpha_u\beta_d\ket{ud}
\\
&\quad
+\alpha_d\beta_u\ket{du}
+\alpha_d\beta_d\ket{dd}.
\label{eq:l6-product-expansion}
\end{align}

\begin{figure}[htbp]
\centering
\includegraphics[width=0.86\textwidth]{lecture_06_figure_04.png}
\caption{Факторизованное двухспиновое состояние, разложенное по составному базису из четырёх состояний.}
\lectureframe{00:37:45}
\end{figure}

\begin{classroomqa}
\audiencequestion{Какого рода произведение векторов появляется в уравнении~\eqref{eq:l6-product-expansion}?}
\lecturerresponse{Рассматриваемое пространство является тензорным произведением двух гильбертовых пространств. Вектор, который сам факторизуется в один вектор из каждого пространства, называется продуктным состоянием или простым тензором. Большинство векторов в пространстве тензорного произведения являются суммами таких произведений и не могут быть записаны как одно произведение.}
\end{classroomqa}

Это ограничение можно увидеть непосредственно по коэффициентам:
\begin{equation}
c_{uu}c_{dd}=c_{ud}c_{du}
\qquad\text{for every product state}.
\label{eq:l6-rank-one-condition}
\end{equation}
Эквивалентно, \(2\times2\)-матрица коэффициентов имеет ранг 1. Подсчёт параметров приводит к тому же выводу. Четыре комплексных коэффициента содержат 8 вещественных параметров; нормировка устраняет 1, а физически несущественная общая фаза устраняет ещё 1, так что для общего чистого состояния двух спинов остаётся 6. Каждому чистому состоянию одного спина нужны 2 вещественных параметра, так что для двух отдельно заданных спинов требуется только 4. Следовательно, состояния-произведения образуют особое 4-параметрическое подмножество 6-параметрического пространства состояний двух спинов. Остальные состояния запутаны.

\section{Локальные наблюдаемые в составном пространстве}

Привычные односпиновые наблюдаемые должны сохраняться и внутри большей системы. Точнее,
\begin{equation}
\sigma_m\equiv\sigma_m\otimes I,
\qquad
\tau_n\equiv I\otimes\sigma_n,
\qquad
m,n\in\{x,y,z\}.
\label{eq:l6-local-operators}
\end{equation}
Сокращённая запись безвредна, как только зафиксировано соглашение о слотах. Используя
\begin{equation}
\begin{aligned}
\sigma_x&=
\begin{pmatrix}0&1\\1&0\end{pmatrix},\\[0.4em]
\sigma_y&=
\begin{pmatrix}0&-i\\i&0\end{pmatrix},\\[0.4em]
\sigma_z&=
\begin{pmatrix}1&0\\0&-1\end{pmatrix},
\end{aligned}
\end{equation}
получаем: операторы первого спина действуют, тогда как вторая метка остаётся пассивной:
\begin{align}
\sigma_x\ket{u\,i}&=\ket{d\,i},
&
\sigma_x\ket{d\,i}&=\ket{u\,i},\\
\sigma_z\ket{u\,i}&=\ket{u\,i},
&
\sigma_z\ket{d\,i}&=-\ket{d\,i}.
\end{align}
Для второго спина пассивной остаётся первая метка:
\begin{align}
\tau_x\ket{a\,u}&=\ket{a\,d},
&
\tau_x\ket{a\,d}&=\ket{a\,u},\\
\tau_z\ket{a\,u}&=\ket{a\,u},
&
\tau_z\ket{a\,d}&=-\ket{a\,d}.
\label{eq:l6-spectator-action}
\end{align}
\(y\)-компонента тоже переворачивает свой активный слот, причём \(\sigma_y\ket{u}=i\ket{d}\) и
\(\sigma_y\ket{d}=-i\ket{u}\).

\begin{figure}[htbp]
\centering
\includegraphics[width=0.86\textwidth]{lecture_06_figure_05.png}
\caption{Таблица действия операторов второго спина: первый индекс остаётся пассивным.}
\lectureframe{00:47:40}
\end{figure}

\begin{classroomqa}
\audiencequestion{Означает ли применение \(\sigma_x\), что произошло измерение \(x\), и почему результат не является суперпозицией?}
\lecturerresponse{Действие оператора само по себе не является измерением. Умножение матрицы \(\sigma_x\) на вектор \(z\)-базиса \(\ket{u}\) даёт \(\ket{d}\). Это не уравнение на собственные значения: собственный вектор сохранял бы своё направление при действии. Действительно, \(\ket{u}\) переходит в ортогональное состояние; это ещё один способ увидеть, почему измерение \(x\) на \(z\)-собственном состоянии максимально случайно. Действие на вектор, не являющийся собственным, обычно представляет собой промежуточный алгебраический шаг, физический смысл которого определяется всей вычисляемой величиной.}
\end{classroomqa}

Операторы, действующие на разные множители, коммутируют. Например,
\begin{equation}
\sigma_x\tau_x\ket{uu}
=\ket{dd}
=\tau_x\sigma_x\ket{uu}.
\end{equation}
То же равенство выполняется на каждом базисном векторе составного пространства и потому на любом состоянии:
\begin{equation}
[\sigma_m,\tau_n]=0
\qquad\text{for all }m,n.
\label{eq:l6-local-commutation}
\end{equation}

\begin{figure}[htbp]
\centering
\includegraphics[width=0.86\textwidth]{lecture_06_figure_06.png}
\caption{Раздельное действие на разные слоты приводит к нулевому коммутатору \([\sigma_x,\tau_x]\).}
\lectureframe{00:56:40}
\end{figure}

\begin{classroomqa}
\audiencequestion{Почему проверка на одном базисном векторе ещё не является доказательством того, что два оператора коммутируют?}
\lecturerresponse{Равенство операторов требует одинакового действия на каждом векторе. Достаточно проверить все векторы некоторого базиса, потому что линейность затем распространяет результат на любую суперпозицию. Для локальных операторов аргумент немедленно проходит на всех 4 базисных состояниях: каждый оператор изменяет другой слот, так что их порядок несущественен.}
\end{classroomqa}

\begin{classroomqa}
\audiencequestion{Существуют ли наблюдаемые, которые действительно принадлежат системе двух спинов?}
\lecturerresponse{Да. Произведения, такие как \(\sigma_x\tau_x\) и \(\sigma_x\tau_y\), а также линейные комбинации таких произведений исследуют совместные свойства. Локальные наблюдаемые описывают каждую подсистему; корреляционные наблюдаемые выявляют, как подсистемы связаны.}
\end{classroomqa}

\begin{classroomqa}
\audiencequestion{Являются ли независимо приготовленные состояния-произведения уже запутанными?}
\lecturerresponse{Нет. В состоянии-произведении у каждого множителя есть свой собственный вектор состояния, и локальные средние значения факторизуются. Запутанность начинается тогда, когда составное состояние нельзя разложить на одно состояние для \(\sigma\) и одно для \(\tau\).}
\end{classroomqa}

\section{Синглетное состояние}

Примером максимально запутанного состояния служит
\begin{equation}
\ket{S}
=\frac{1}{\sqrt2}\bigl(\ket{ud}-\ket{du}\bigr).
\label{eq:l6-singlet}
\end{equation}
Эти два члена уже показывают идеальную антикорреляцию вдоль \(z\): всякий раз, когда один спин направлен вверх, другой направлен вниз. Более глубокий момент состоит в том, что ни одному из спинов нельзя приписать собственное чистое состояние.

\begin{figure}[htbp]
\centering
\includegraphics[width=0.86\textwidth]{lecture_06_figure_07.png}
\caption{Антисимметричное двухспиновое состояние, которое будет отождествлено с синглетом.}
\lectureframe{01:09:25}
\end{figure}

Для чистого состояния одного спина всегда существует единичный вектор \(\bm n\), для которого
\begin{equation}
(\bm\sigma\cdot\bm n)\ket{\psi}=\ket{\psi}.
\end{equation}
Эквивалентно, вектор Блоха
\begin{equation}
\bm r
=\bigl(\expect{\sigma_x},\expect{\sigma_y},\expect{\sigma_z}\bigr)
\end{equation}
для чистого односпинового состояния имеет длину 1. Поэтому все три его компоненты не могут одновременно обращаться в нуль. Для данного состояния ориентация детектора вдоль \(\bm r\) даёт \(+1\) с достоверностью. Знакомые примеры: \(z\) для \(\ket{u}\), \(x\) для \((\ket{u}+\ket{d})/\sqrt2\) и одно из \(y\)-направлений для \((\ket{u}+i\ket{d})/\sqrt2\).

Теперь вычислим локальное среднее в синглете:
\begin{align}
\expect{\sigma_z}_S
&=\frac12
(\bra{ud}-\bra{du})\sigma_z
(\ket{ud}-\ket{du})\\
&=\frac12
(\bra{ud}-\bra{du})
(\ket{ud}+\ket{du})\\
&=0.
\end{align}
Для \(\sigma_x\),
\begin{equation}
\sigma_x\ket{S}
=\frac{\ket{dd}-\ket{uu}}{\sqrt2},
\end{equation}
что ортогонально \(\ket{S}\), так что \(\expect{\sigma_x}_S=0\). Фазы, вносимые \(\sigma_y\), не меняют этой ортогональности, поэтому
\begin{equation}
\expect{\sigma_x}_S
=\expect{\sigma_y}_S
=\expect{\sigma_z}_S=0.
\label{eq:l6-zero-local-sigma}
\end{equation}
По симметрии,
\begin{equation}
\expect{\tau_x}_S
=\expect{\tau_y}_S
=\expect{\tau_z}_S=0.
\label{eq:l6-zero-local-tau}
\end{equation}

\begin{figure}[htbp]
\centering
\includegraphics[width=0.86\textwidth]{lecture_06_figure_08.png}
\caption{На доске записаны нулевые локальные средние значения спина и переворот по \(y\)-компоненте, использованный для завершения проверки.}
\lectureframe{01:14:20}
\end{figure}

Всё состояние двух спинов чисто и полностью известно, но каждая его часть локально неполяризована. Поскольку ни одно чистое односпиновое состояние не имеет нулевого вектора Блоха, \(\ket{S}\) не может быть состоянием-произведением.

\begin{classroomqa}
\audiencequestion{Имело бы обычное состояние-произведение направление с ненулевой локальной поляризацией?}
\lecturerresponse{Да. Доказательство для одного спина применимо к соответствующему множителю, тогда как другой множитель лишь остаётся пассивным. Поэтому каждый множитель чистого состояния-произведения имеет единичный вектор Блоха. Синглет не обладает этим свойством состояния-произведения для обоих спинов.}
\end{classroomqa}

\begin{classroomqa}
\audiencequestion{Как можно приготовить два спина в синглете?}
\lecturerresponse{Связь, благоприятствующая антивыравниванию, может сделать синглет состоянием наименьшей энергии. Приведите спины во взаимодействие и дайте паре терять энергию в окружающую среду; тогда она может релаксировать в синглет. Это утверждение зависит от взаимодействия. Для изотропного антиферромагнитного гамильтониана \(H_{\rm int}=J\,\bm\sigma\cdot\bm\tau\) при \(J>0\) синглет лежит ниже трёх триплетных состояний.}
\end{classroomqa}

\begin{classroomqa}
\audiencequestion{Сколько времени занимает такая релаксация?}
\lecturerresponse{Характерный масштаб времени зависит от расстояния между спинами, силы их взаимной связи, их магнитных моментов и их связи с электромагнитным окружением, которое уносит энергию. Более сильная связь обычно означает более быструю релаксацию.}
\end{classroomqa}

\begin{classroomqa}
\audiencequestion{Почему антисимметричная комбинация может иметь меньшую энергию, чем симметричная?}
\lecturerresponse{Для двух частиц со спином \(\tfrac12\) антисимметричное состояние имеет полный спин 0, тогда как три симметричных состояния образуют триплет с полным спином 1. Взаимодействие, пропорциональное \(+\bm\sigma\cdot\bm\tau\), приписывает синглету собственное значение \(-3\), а триплету — \(+1\), так что положительная связь располагает синглет ниже. Другие гамильтонианы не обязаны упорядочивать уровни таким же образом.}
\end{classroomqa}

\subsection{Одно и то же состояние в любом спиновом базисе}

Запутанность не устраняется выбором более удобного локального базиса. Независимой от базиса мерой, которая будет введена позже, служит энтропия запутанности. Синглет проявляет ещё более сильную симметрию. Подстановка
\begin{equation}
\ket{u}=\frac{\ket{R}+\ket{L}}{\sqrt2},
\qquad
\ket{d}=\frac{\ket{R}-\ket{L}}{\sqrt2},
\end{equation}
в уравнение~\eqref{eq:l6-singlet} даёт
\begin{equation}
\ket{S}
=\frac{\ket{LR}-\ket{RL}}{\sqrt2}.
\label{eq:l6-singlet-x-basis}
\end{equation}
В любом общем ортонормированном спиновом базисе \(\{\ket{\rm in},\ket{\rm out}\}\),
\begin{equation}
\ket{S}
=\frac{\ket{{\rm in}\,{\rm out}}-\ket{{\rm out}\,{\rm in}}}{\sqrt2},
\label{eq:l6-singlet-any-basis}
\end{equation}
с точностью до физически несущественной общей фазы. Три триплетных состояния при общем повороте перемешиваются между собой; одномерный синглетный сектор может лишь воспроизводить сам себя.

\begin{figure}[htbp]
\centering
\includegraphics[width=0.86\textwidth]{lecture_06_figure_09.png}
\caption{Антисимметричная структура сохраняется при переходе от базиса вверх-вниз к базису вправо-влево и к другим общим спиновым базисам.}
\lectureframe{01:21:00}
\end{figure}

\begin{classroomqa}
\audiencequestion{Является ли запутанность всего лишь артефактом базиса, в котором записано состояние?}
\lecturerresponse{Нет. Локальные замены базиса изменяют коэффициенты, но не то, факторизуется ли состояние, и энтропия запутанности количественно выражает этот инвариантный факт. Для синглета общие повороты делают эту инвариантность особенно наглядной: антисимметричная форма воспроизводится в каждом спиновом базисе.}
\end{classroomqa}

\section{Корреляционные наблюдаемые}

Локальные средние обращаются в нуль, так что информация в синглете должна содержаться в совместных наблюдаемых. Начнём с
\begin{equation}
\sigma_z\tau_z\ket{ud}=-\ket{ud},
\qquad
\sigma_z\tau_z\ket{du}=-\ket{du}.
\end{equation}
Следовательно,
\begin{equation}
\sigma_z\tau_z\ket{S}=-\ket{S},
\qquad
\expect{\sigma_z\tau_z}_S=-1.
\label{eq:l6-zz-correlation}
\end{equation}
Результат — не просто среднее: синглет является собственным состоянием этого произведения, поэтому измерения \(z\) вдоль одной и той же оси всегда дают противоположные знаки.

Оба \(x\)-оператора переворачивают свои слоты:
\begin{equation}
\sigma_x\tau_x\ket{ud}=\ket{du},
\qquad
\sigma_x\tau_x\ket{du}=\ket{ud}.
\end{equation}
Перестановка двух членов меняет знак антисимметричного состояния:
\begin{equation}
\sigma_x\tau_x\ket{S}=-\ket{S},
\qquad
\expect{\sigma_x\tau_x}_S=-1.
\end{equation}
Две фазы от \(y\)-операторов сокращаются точно так же, давая
\begin{equation}
\expect{\sigma_y\tau_y}_S=-1.
\end{equation}
Напротив, для различных декартовых компонент среднее равно 0. Например,
\begin{equation}
\sigma_x\tau_z\ket{S}
=-\frac{\ket{dd}+\ket{uu}}{\sqrt2},
\qquad
\expect{\sigma_x\tau_z}_S=0.
\end{equation}
Все эти результаты суммируются формулой
\begin{equation}
\expect{\sigma_m\tau_n}_S=-\delta_{mn},
\qquad
\expect{(\bm\sigma\cdot\bm a)(\bm\tau\cdot\bm b)}_S
=-\bm a\cdot\bm b.
\label{eq:l6-singlet-correlation-tensor}
\end{equation}

\begin{figure}[htbp]
\centering
\includegraphics[width=0.86\textwidth]{lecture_06_figure_10.png}
\caption{В вычислении корреляций получаются \(\expect{\sigma_x\tau_x}=-1\) и результат для смешанных осей \(\expect{\sigma_x\tau_z}=0\).}
\lectureframe{01:32:20}
\end{figure}

\begin{classroomqa}
\audiencequestion{Что физически означает произведение двух наблюдаемых? Выполняется ли одно измерение после другого?}
\lecturerresponse{Вообще говоря, эрмитово произведение — это другая наблюдаемая в составном пространстве. Здесь каждый \(\sigma_m\) коммутирует с каждым \(\tau_n\), поэтому обе локальные компоненты можно измерить одновременно и перемножить их численные результаты. Оператор произведения кодирует этот совместный результат; никакого временного упорядочения не требуется.}
\end{classroomqa}

Для состояния-произведения совместные средние факторизуются:
\begin{equation}
\expect{\sigma_m\tau_n}_{\rm prod}
=\expect{\sigma_m}_{\sigma}\expect{\tau_n}_{\tau}.
\end{equation}
Синглет максимально нарушает такую факторизацию. Каждый множитель справа был бы равен нулю, тогда как произведение слева для одной и той же оси равно \(-1\). По отдельности спины полностью случайны; совместно же одинаковые компоненты заведомо противоположны.

\begin{classroomqa}
\audiencequestion{Всякое ли взаимодействие между двумя спинами порождает запутанность?}
\lecturerresponse{Типичное взаимодействие может запутать первоначально факторизованное состояние, и многие обычные взаимодействия именно так и делают. Однако это не является буквально неизбежным: специально выбранное собственное состояние-произведение может лишь приобрести фазу, а чисто локальный гамильтониан не может запутать подсистемы. Существенно именно нетривиальное взаимодействие, действующее на состояние, которое оно не сохраняет в виде произведения.}
\end{classroomqa}

\section{От корреляции к вопросу Белла}

Сама по себе совершенная антикорреляция не парадоксальна. Представьте, что в карман кладут пенни и дайм, одну монету, не глядя, отдают Алисе, а другую Бобу и разводят их далеко друг от друга. Когда Алиса видит, какая монета у неё, она сразу знает, какая у Боба. Корреляция была установлена, когда монеты находились вместе; узнавание одного результата не посылает никакого сигнала в другое место.

Синглет имеет ту же элементарную особенность, но добавляет решающую свободу: Алиса и Боб могут выбирать любые направления измерения, и квантовая механика предсказывает непрерывно меняющуюся корреляцию \(-\bm a\cdot\bm b\). Теорема Белла ставит вопрос, могут ли все эти корреляции возникать из локальных классических инструкций, которые частицы несут с собой заранее.

\begin{classroomqa}
\audiencequestion{Именно это состояние используется в теореме Белла?}
\lecturerresponse{Да. В рассуждениях типа Белла используют запутанные пары, такие как синглет, а не определённое состояние-произведение \(\ket{ud}\). Это различие важно, потому что состояние-произведение даёт раздельные локальные описания, тогда как синглет несёт единую неразделимую структуру корреляций в любом общем базисе.}
\end{classroomqa}

\begin{classroomqa}
\audiencequestion{Почему антикорреляцию называют статистической, если результаты для одной и той же оси всегда противоположны?}
\lecturerresponse{Каждый локальный результат случаен: Алиса не может предсказать, получит ли она \(+1\) или \(-1\). Условная связь точна: как только известен один результат при одинаковых осях, другой фиксирован как противоположный. Корреляция относится к совместному распределению, а не к неопределённости произведения при одинаковых осях.}
\end{classroomqa}

Удобный способ заострить вопрос — рассмотреть моделирование. Один классический компьютер может имитировать любой эксперимент с одним спином: хранить его состояние, эволюционировать его по уравнению Schr\"odinger, вычислять требуемые вероятности и с помощью генератора случайных чисел выбирать результат. Можно воспроизвести и повторные измерения, и магнитную прецессию.

Теперь разделите запутанную пару между двумя классическими компьютерами. Позвольте компьютерам взаимодействовать при подготовке пары, затем разъедините их и независимо выберите направления измерения в удалённых местах. Никакая пара разделённых локальных классических программ, даже с общей случайностью, заранее зафиксированной, не может воспроизвести квантовые корреляции при любом выборе направлений. Глобальное классическое моделирование возможно, но для этого нужно сохранять доступ к совместному состоянию или иначе координировать обе стороны так, что модель уже не будет локальной моделью со скрытыми переменными.

Эта несостоятельность локального классического учёта не даёт канала связи быстрее света. Локальные результаты Алисы остаются случайными, поэтому она не может управлять сообщением, которое Боб смог бы прочитать без последующего сравнения их записей. Терминология вокруг нелокальности различается, но операционное различие ясно: корреляции Белла исключают локальные классические объяснения, сохраняя при этом отсутствие сигнализации.

Тем самым построение составного состояния породило нечто действительно новое. Тензорные произведения локально сохраняют прежнюю физику одного спина, однако их типичные векторы нельзя разделить на отдельные состояния. В синглете вся локальная поляризация исчезает, тогда как совместные компоненты спина несут совершенную, не зависящую от базиса антикорреляцию. Следующий шаг — превратить эту структуру в количественный тест Белла для классического рассуждения.
